% =========================================================
% III. PROBLEM FORMULATION
% =========================================================
\section{Phát biểu bài toán}
\label{sec:problem_formulation}

Chúng tôi xét bài toán hỏi đáp pháp luật nhiều bước, trong đó một kết luận có
thể yêu cầu liên kết các điều kiện luật định, sự kiện trung gian và hệ quả pháp
lý. Tri thức pháp luật được biểu diễn bằng đồ thị có hướng $G=(V,E)$ và tập quy
tắc $\mathcal{R}$. Mỗi nút biểu diễn một sự kiện pháp lý đã được chuẩn hóa, còn
mỗi cạnh tóm tắt một hoặc nhiều quy tắc điều kiện--hệ quả. Đồ thị hỗ trợ truy
xuất và phân tích tô-pô, trong khi $\mathcal{R}$ xác định ngữ nghĩa suy luận.

Với câu hỏi $q$, hệ thống xây dựng các đường suy luận ứng viên
\begin{equation}
\mathcal{P}(q)=\{P_1,\ldots,P_K\},\qquad
P_i=v_0\xrightarrow{r_1}\cdots\xrightarrow{r_h}v_h,
\end{equation}
và chọn đường chính
\begin{equation}
P^*=\arg\max_{P_i\in\mathcal{P}(q)}S(P_i,q),
\end{equation}
trong đó $S$ kết hợp độ phù hợp ngữ nghĩa và cấu trúc.

\subsection{Ngữ nghĩa ba giá trị dựa trên đại số tuyến tính}

Mỗi sự kiện được gắn với một biến
\begin{equation}
X_i\in\{1,0,\bot\},
\end{equation}
lần lượt biểu diễn \textsc{true}, \textsc{false} và \textsc{unknown}. Trạng
thái unknown thể hiện hỗ trợ bị thiếu hoặc xung đột, không phải một giá trị số
nằm giữa 0 và 1. Formulation này dựa trên ngữ nghĩa điểm bất động ba giá trị
\cite{kleene1952metamathematics,fitting1985kripke} và các mở rộng trong không
gian vector \cite{sato2020threevalued}.

Trạng thái sự kiện được mã hóa bằng hai kênh hỗ trợ:
\begin{equation}
(t_i,f_i)=
\begin{cases}
(1,0), & X_i=1,\\
(0,1), & X_i=0,\\
(0,0), & X_i=\bot.
\end{cases}
\label{eq:two_channel_state}
\end{equation}

Đặt $n=|V|$ và $p=|\mathcal{R}|$. Chúng tôi định nghĩa các ma trận incidence
cho thân quy tắc
$\mathbf{B}^{+},\mathbf{B}^{-}\in\{0,1\}^{p\times n}$ và các ma trận incidence
cho đầu quy tắc
$\mathbf{H}^{+},\mathbf{H}^{-}\in\{0,1\}^{n\times p}$. Với
$\mathbf{k}=(\mathbf{B}^{+}+\mathbf{B}^{-})\mathbf{1}$, các thân quy tắc và mức
kích hoạt được tính bởi
\begin{equation}
\mathbf{b}
=
\mathbb{I}
[\mathbf{B}^{+}\mathbf{t}
+\mathbf{B}^{-}\mathbf{f}
=\mathbf{k}],
\qquad
\mathbf{a}
=
\mathbf{b}\odot(\mathbf{1}-\mathbf{e}),
\label{eq:rule_activation}
\end{equation}
trong đó $\mathbf{e}$ biểu diễn các ngoại lệ đang hoạt động. Điều kiện bằng
$\mathbf{k}$ yêu cầu mọi literal trong thân quy tắc đều được thỏa mãn, qua đó
biểu diễn phép hội trong từng quy tắc.

Hỗ trợ cho các đầu quy tắc được tổng hợp như sau:
\begin{equation}
\mathbf{d}^{+}
=
\mathbb{I}[\mathbf{H}^{+}\mathbf{a}\geq1],
\qquad
\mathbf{d}^{-}
=
\mathbb{I}[\mathbf{H}^{-}\mathbf{a}\geq1].
\label{eq:head_support}
\end{equation}
Phép đặt ngưỡng biểu diễn phép tuyển giữa các quy tắc được kích hoạt độc lập có
cùng hệ quả. Do đó, $(A\land B)\rightarrow Y$ được mã hóa thành một quy tắc có
hai phần tử trong thân, còn $(A\lor B)\rightarrow Y$ được chuẩn hóa thành
$A\rightarrow Y$ và $B\rightarrow Y$. Biểu diễn incidence này có liên hệ với
lập trình logic trong không gian ma trận và tensor
\cite{sakama2021tensor}.

Trạng thái factual được xác định bằng phép lặp tới điểm bất động:
\begin{equation}
\mathbf{x}^{F}
=
\operatorname{Fix}
(\Phi_{\mathcal{R}},\mathbf{x}^{(0)}).
\label{eq:factual_fixed_point}
\end{equation}
Chỉ có hỗ trợ dương hoặc chỉ có hỗ trợ âm lần lượt tạo trạng thái true hoặc
false; không có hỗ trợ hoặc có hỗ trợ xung đột tạo trạng thái unknown. Runtime
đánh giá các quy tắc thưa bằng phép suy luận ký hiệu thay vì thực hiện phép nhân
ma trận dày. Mặc dù engine hỗ trợ thân quy tắc có nhiều literal, dữ liệu hiện
tại chỉ chứa một tiền đề dương trong mỗi quy tắc được import. Vì vậy, khả năng
xử lý phép hội chưa được đánh giá thực nghiệm trong nghiên cứu này.

\subsection{Chẩn đoán tô-pô và can thiệp cấu trúc}

Gọi $\mathcal{M}(P^*)=\{v_1,\ldots,v_{h-1}\}$ là tập sự kiện trung gian của
đường chính và $m\in\mathcal{M}(P^*)$ là một sự kiện trung gian. Phép chẩn đoán
xóa nút tạo
\begin{equation}
G^{-m}=G[V\setminus\{m\}]
\end{equation}
và kiểm tra
\begin{equation}
\rho_{\mathrm{top}}(s,Y\mid m)
=
\mathbb{I}
[Y\text{ có thể đạt được từ }s\text{ trong }G^{-m}].
\label{eq:topological_ablation}
\end{equation}
Phép toán này chỉ đo tính đạt được tô-pô có giới hạn và không được sử dụng làm
định nghĩa phản thực của LegalSCM.

Với câu hỏi yêu cầu loại bỏ sự kiện trung gian một cách tường minh, LegalSCM
thay vào đó áp dụng can thiệp cấu trúc \cite{pearl2009causality}
\begin{equation}
do(X_m=0),\qquad 0\equiv\mathrm{FALSE}.
\end{equation}
Can thiệp cố định $X_m$ ở trạng thái false và vô hiệu hóa mọi quy tắc xác định
biến này:
\begin{equation}
\mathcal{R}^{do(X_m=0)}
=
\{r\in\mathcal{R}\mid
\operatorname{var}(\operatorname{head}(r))\neq m\}.
\end{equation}
Sau đó, trạng thái phản thực được suy luận lại:
\begin{equation}
\mathbf{x}^{CF(m)}
=
\operatorname{Fix}\!\left(
\Phi_{\mathcal{R}^{do(X_m=0)}},
\mathbf{x}^{(0)}[X_m\leftarrow0]
\right).
\label{eq:cf_fixed_point}
\end{equation}

Sự kiện trung gian $m$ là cần thiết đối với đích $Y$ khi $x_Y^{F}=1$ và
$x_Y^{CF(m)}=0$; không cần thiết khi cả hai giá trị đều bằng $1$; và chưa xác
định trong các trường hợp còn lại. Đây là can thiệp cấu trúc tương đối với mô
hình đã mã hóa, không có bước abduction đối với biến ngoại sinh, và không được
đồng nhất với thao tác xóa $m$ khỏi $G$.

Cuối cùng, hệ thống dự đoán
\begin{equation}
z\in\{
\textsc{supported},
\textsc{reject-direct-claim},
\textsc{uncertain}
\}
\end{equation}
và sinh đầu ra
\begin{equation}
f:(q,G,\mathcal{R})
\rightarrow(P^*,z,\mathcal{E}^*,a,C),
\end{equation}
trong đó $\mathcal{E}^*$ là bằng chứng đã được xác minh, $a$ là câu trả lời và
$C$ là tập căn cứ pháp lý được trích dẫn.
